\section{The decomposition hypergraph functor \(\mathbf{D}\)}\label{sec:functor}

We now construct a contravariant functor
\[
\mathbf{D} : \mathcal{A}^{\mathrm{op}} \longrightarrow \mathcal{H}.
\]

\subsection{On objects}
For an observable \(o \in \mathcal{A}\), define the set
\[
\mathcal{L}(o) = \{\, x \in \mathcal{O} \mid \text{there exists a directed path } o \rightsquigarrow x \text{ in } \mathcal{G} \,\}.
\]
In words, \(\mathcal{L}(o)\) consists of all observables into which \(o\) can be decomposed (in one or more steps); i.e.\ the parts of \(o\) and their parts, recursively.
The \emph{decomposition hypergraph} of \(o\) is the subgraph of \(\mathcal{G}\) induced by \(\mathcal{L}(o)\):
\[
\mathbf{D}(o) = (\mathcal{L}(o), E_o), \qquad
E_o = \{\, (x, p) \in \mathcal{L}(o) \times \mathcal{L}(o) \mid (x,p) \in E(\mathcal{G}) \,\}.
\]

\subsection{On morphisms}
Let \(\phi : o \to o'\) be a morphism in \(\mathcal{A}\), i.e.\ a decomposition path from \(o\) to \(o'\).
In \(\mathcal{A}^{\mathrm{op}}\) this corresponds to a morphism \(\phi^{\mathrm{op}} : o' \to o\).
We must define a graph homomorphism \(\mathbf{D}(\phi^{\mathrm{op}}) : \mathbf{D}(o') \to \mathbf{D}(o)\).

If there is a path from \(o'\) to \(x\), then by concatenating with the reverse of the path \(\phi\) (which goes from \(o\) to \(o'\) in \(\mathcal{A}\)) we obtain a path from \(o\) to \(x\). Hence \(\mathcal{L}(o') \subseteq \mathcal{L}(o)\).
Moreover, any edge between vertices of \(\mathcal{L}(o')\) is still present in \(\mathbf{D}(o)\) because the edge set of \(\mathbf{D}(o)\) contains all original edges among vertices that belong to \(\mathcal{L}(o)\).
Therefore the inclusion map \(\iota : \mathcal{L}(o') \hookrightarrow \mathcal{L}(o)\) is a graph homomorphism.
We set
\[
\mathbf{D}(\phi^{\mathrm{op}}) = \iota .
\]

\subsection{Functoriality}
For the identity \(\mathrm{id}_o\) in \(\mathcal{A}\), we have \(\mathcal{L}(o) \subseteq \mathcal{L}(o)\) trivially and the inclusion is the identity on \(\mathbf{D}(o)\).
For a composition \(\psi \circ \phi : o \to o''\) in \(\mathcal{A}\), the chain of inclusions \(\mathcal{L}(o'') \hookrightarrow \mathcal{L}(o') \hookrightarrow \mathcal{L}(o)\) coincides with the inclusion \(\mathcal{L}(o'') \hookrightarrow \mathcal{L}(o)\).
Thus \(\mathbf{D}\) is a contravariant functor.